\section{Related Work}

\paratitle{Hallucination Detection.}
Hallucination Detection in LLMs is a pivotal concern due to its role in identifying inaccuracies or falsehoods within model responses. Most existing methods for hallucination detection are implemented based on powerful LLMs~\cite{luo2024hallucination}. One category of these methods relies on the internal knowledge and consistency of LLMs to detect hallucinations, by breaking down the detection process~\cite{dhuliawala2023chain} or comparing multiple responses to the same query~\cite{manakul2023selfcheckgpt}. 
However, these approaches are inherently limited by the knowledge boundaries of LLMs. Another category of methods leverages external tools for hallucination detection~\cite{chern2023factool, wei2024long}. 
While tool utilization can compensate for the knowledge limitations of LLMs, existing methods often require manually designing specific tools tailored to particular text types. In contrast, HaluAgent is equipped with a versatile and expandable toolbox, enabling small language models to autonomously select appropriate tools for fine-grained detection.

\paratitle{Agent Tuning.}
Recently, the surprising planning and reasoning capabilities of LLMs have inspired research into their application as agents for specific tasks. 
Previous research~\cite{yao2023react, nakano2021webgpt, singh2023progprompt, yang2023mm} has predominantly relied on prompting to use LLMs as agents, but this method demands advanced instruction-following capabilities that open-source LLMs typically do not match with API-based LLMs.
To break this restriction, AgentTuning~\cite{zeng2023agenttuning} first fine-tunes open-source LLMs on agent interaction trajectories generated by powerful LLMs. Moreover, Agent-FLAN~\cite{chen2024agent} goes a step further by decomposing and redesigning the training corpus.
Besides general agents, KG-Agent~\cite{KGAgent} fine-tunes LLaMA-7B model to achieve an agent specialized in reasoning over knowledge graphs. Similarly, HaluAgent is fine-tuned on trajectory data of hallucination detection tasks, thereby enhancing the detection capabilities of open-source LLMs.